\documentclass[11pt]{article}

\usepackage[T1]{fontenc}

\usepackage{lmodern}

\usepackage[margin=22mm]{geometry}

\usepackage{booktabs,tabularx,array}

\usepackage{hyperref}

\usepackage{microtype}

\hypersetup{colorlinks=true,urlcolor=blue}

\setlength{\parskip}{0.55em}

\setlength{\parindent}{0pt}

\title{A provenance-preserving readiness audit of model-built NNC2215 structures for a paracetamol glucuronide study\\[0.4em]\large Public-data reanalysis, coordinate mapping, and force-field review gates}

\author{Author list and affiliations pending contributor confirmation}

\date{26 October 2026}

\begin{document}

\maketitle

\begin{center}\fbox{\parbox{0.92\textwidth}{\small Internal scientific manuscript draft. Author list, affiliations, contributions, funding, and disclosures require contributor confirmation. This is not a validated drug-discovery result and is not ready for journal submission.}}\end{center}

\begin{abstract}

NNC2215 is a published glucose-sensitive insulin conjugate. A 2026 study reported interaction with paracetamol glucuronide (PaG), motivating a review of structural reproducibility and model readiness while keeping pharmacological claims tied to experiments. We audited published source data, publisher-supplied model-built open and closed structures, a source-derived heavy-atom graph, conditional charge-state bookkeeping, and limited fragment-parameterization pilots. Reanalysis of three independent Figure 3d experiments gave paired mean fold changes of 3.1576 (sample SD 0.2786) from 3 to 20 mmol/L glucose and 12.4720 (SD 1.6985) from 0 to 20 mmol/L, consistent with published rounded values. The 507-heavy-atom graph comprises 398 insulin atoms and 109 conjugate atoms and maps to both model-built candidates. One B1 mapping and 16 B29 mappings per state remain in the ledger; the B29 alternatives are not resolved by coordinates. A group-level pH 7.4 audit enumerates 144 group-charge combinations without assigning populations. Fragment parameterization and OpenFF label checks do not constitute a validated hybrid force field. No full OpenMM target system, target-binding molecular dynamics, affinity estimate, or clinical result was produced. The current decision is no-go for target-binding MD pending independent chemical and parameter review. This is a reproducibility and readiness assessment of an existing published candidate, not discovery of a new drug.

\end{abstract}

\textbf{Keywords:} NNC2215; paracetamol glucuronide; insulin conjugate; coordinate mapping; force-field readiness; reproducibility

\bigskip

\section{Introduction}

NNC2215 is an insulin conjugate with a glucose-responsive switch formed by a glucose-binding macrocycle linked at LysB29 and a glucoside linked at PheB1. The 2024 report described glucose-dependent receptor binding and animal studies [1]. A 2026 study reported stronger binding of NNC2215 to paracetamol glucuronide (PaG) than to D-glucose in vitro and increased insulin-receptor affinity at 70 micromol/L PaG. Its authors concluded NNC2215 was not considered a viable therapeutic candidate because of potential interaction-related hypoglycemia risk [2]. This work does not revisit those experimental conclusions.

The computational question here is narrower: what public evidence is sufficient to prepare a reproducible, chemically reviewed model for a future, explicitly labelled proxy study? Public open and closed coordinates are model-built representations. The 2026 supporting-information record describes a 500 ns Desmond simulation of a related glucose-bound macrocycle analogue, but the checked public record does not list its in-house starting coordinates, system, trajectory, or selected snapshot as machine-readable files [2,4]. A simulation using the public NNC2215 candidates would therefore be a separate proxy analysis, not an exact rerun.

We report an auditable readiness assessment: arithmetic reanalysis of published measurements, source-derived heavy-atom mapping, explicit retention of unresolved alternatives, and a review of conditions required before system construction or dynamics.

\section{Materials and methods}

\subsection{Evidence corpus and published-data arithmetic}

The evidence corpus comprised the 2024 NNC2215 article and public structural-model record [1,3], the 2026 PaG article and supporting-information record [2,4], and project source files and audit outputs at repository baseline commit 9675159eefc60bc4a8a12661bce208b99d342246 [5]. The published Figure 3d workbook contains three independent experiments. We checked all 30 source cells, calculated 20-to-3 and 20-to-0 mmol/L ratios within each experiment, then report the mean paired ratio and sample standard deviation. Published PaG binding summaries were transcribed as reported; raw binding curves and individual animal observations were not refit.

\subsection{Source graph and coordinate mapping}

A source-derived constitutional graph was reconciled against the insulin sequence, published modification descriptions, and atom identities in the open and closed publisher model files. The graph contains 507 heavy atoms: 398 insulin atoms and 109 conjugate atoms. The coordinate ledger retains alternative exact element-adjacency mappings, records all 527 source-edge distances, and does not use PDB CONECT records to infer bond order, formal charge, or unresolved hydrogen ownership. Broad C/N/O distance screens cover 518 edges; the nine sulfur-containing edges are retained as observations but left unscreened because no thresholds were defined. The separate 12-heavy-atom BGL record in the open candidate was tracked as unlinked to the conjugate graph.

The attachment-boundary audit cross-checked B1, B29, and roof/linker joins in both candidates and preserved one exact B1 mapping and 16 exact B29 mappings per state. It compared the PheB1 N-terminal hydrogen records by coordinate and local geometry. Geometry is treated as compatibility evidence only; it cannot establish atom identity or chemical ownership.

\subsection{Protonation and parameter-feasibility audits}

The pH 7.4 charge-state ledger enumerates combinations of three macrocycle carboxyl groups, the A-chain N-terminus, and two histidines. It records formal-charge bookkeeping and hydrogen counts only; it assigns no pKa values or state probabilities. Amber14 residue-template inventory, OpenFF Sage 2.3.0 label coverage, and capped B1/B29 GAFF2/AM1-BCC fragment pilots were used as feasibility diagnostics [6,7]. These checks do not validate a complete protein-conjugate force field, atom-localized charges, cross-boundary terms, or transferability.

\subsection{Readiness decision rule}

Target-binding molecular dynamics was considered blocked unless an independent reviewer approves the atom-mapped graph, stereochemistry, selected microstate, formal charge, modified-residue boundaries, charges, bonded terms, impropers, and protein-force-field compatibility. A subsequent OpenMM system would require complete topology and parameter coverage, correct atom and charge totals, acceptable geometry and clashes, and finite energies before any restrained minimization or stability pilot.

\section{Results}

\subsection{Published experimental context}

The source-cell reanalysis yielded a mean within-experiment fold change of 3.1576 (SD 0.2786, n = 3) from 3 to 20 mmol/L glucose and 12.4720 (SD 1.6985, n = 3) from 0 to 20 mmol/L, matching the 2024 paper rounded to 3.2-fold and 12.5-fold [1]. These are arithmetic checks on published data, not new experiments.

The 2026 article reports mean dissociation constants of 1.91 mmol/L for D-glucose and 0.17 mmol/L for PaG; their ratio is approximately 11.2 when calculated from the rounded means. At zero glucose, the reported hIR-A relative-affinity summary changes from 0.9\% without PaG to 5.4\% with 70 micromol/L PaG [2]. These published summaries were not independently refit. The 2026 authors state that NNC2215 was not considered viable for therapeutic use in light of the potential interaction risk [2].

\subsection{Coordinate coverage and unresolved assignments}

The source-derived 507-heavy-atom graph maps to both corrected, model-built candidates. Each state retains one B1 assignment and 16 B29 assignments; all 527 graph-edge distances are recorded for each mapping. Broad distance screens cover 518 C/N/O edges; the nine sulfur-containing edges remain unscreened. This closes heavy-atom coordinate coverage under the declared graph and mapping rules, not chemical identity for every ambiguous atom. The free BGL record in the open model remains separate and unlinked.

For the PheB1 N-terminal hydrogen, open-model H2 is 1.012274 angstrom from PheB1 N and has an explicit N CONECT edge. Closed-model PHX H2A is 1.006376 angstrom from that N but has no N CONECT edge. Similar local geometry does not decide whether the closed hydrogen belongs to PheB1 N. The PDBs were not edited to resolve this ambiguity.

\begin{center}\begin{minipage}{\textwidth}
\small\textbf{Table 1. Coordinate and parameter readiness evidence}\par\smallskip
\centering\begin{tabularx}{\textwidth}{>{\raggedright\arraybackslash}p{0.19\textwidth} X X}\toprule

Area & Observed evidence & Open review question \\ \midrule

Graph and coordinates & 507 heavy atoms: 398 insulin plus 109 conjugate; one B1 and 16 B29 mappings per state; 527 graph-edge distances retained. & Review source interpretation, stereochemistry, candidate edits, and B29 alternatives. \\

Group-charge combinations & 144 combinations; net formal-charge range -8 to -2 in the declared enumeration. & No pKa or population evidence; select and justify atom-level states and total system charge. \\

Protein templates & Amber template-name inventory covers 16 of 17 observed residue labels; HIE and CYX are context-dependent candidates. & Validate termini, disulfides, modified residues, and cross-boundary terms in one topology. \\

B1 fragment pilot & 37 heavy atoms; GAFF2/AM1-BCC partial-charge sum +0.003002 e for the declared capped fragment. & Truncation cap and local fragment prevent interpreting this as final conjugate charge assignment. \\

B29 fragment pilot & 92 heavy atoms in declared -3 state; eight label variants; charge sums -2.997001 to -2.996005 e; two improper records have penalty score 41.9. & Review mapping, cap, impropers, charge localization, and cross-boundary terms; no variant is selected. \\

Full-system readiness & OpenFF topology-derived label counts were inventoried; no validated full-conjugate force field exists. & No full OpenMM target system, target MD, affinity estimate, or clinical result. \\

\bottomrule\end{tabularx}
\end{minipage}\end{center}

\subsection{Parameterization and system gate}

Amber template and OpenFF label inventories indicate that standard protein terms can be located or typed, but custom PheB1, LysB29, and linker/macrocycle terms remain decisive. The B1 and B29 pilots are capped fragments with declared local assumptions. Across the 16 B29 atom mappings, calculated partial-charge sums range from -2.997001 to -2.996005 e. Within each mapping, the eight carboxylate resonance label variants have identical charge sums. No charge correction or averaging was applied, and no mapping was selected as authoritative. The penalty score of 41.9 is a parameter-generation diagnostic requiring expert inspection, not a physical energy.

The research record contains no complete force-field assignment and no OpenMM target system. Consequently, no target minimization, equilibration, PaG pose test, target-binding MD, binding free energy, or target-specific GPU throughput measurement has been produced.

\section{Discussion}

This audit reduces uncertainty in coordinate coverage while keeping chemical uncertainty visible. A single heavy-atom graph can map to multiple coordinate assignments, and geometric agreement cannot determine atom identity, bond order, protonation, charge localization, or force-field suitability. The retained B29 alternatives and PheB1 hydrogen discrepancy are substantive review items.

The smallest defensible next step is independent review by a qualified computational chemist of the atom-mapped conjugate, stereochemistry, B1/B29 boundaries, hydrogen assignment, selected microstate, full formal charge, and candidate fragment parameters. If an explicit state is approved, the next computational gate is one complete OpenMM topology check for atom and parameter coverage, charge total, clashes, and finite energy. Only after that passes should restrained minimization and a short stability check be considered. A later apo/glucose/PaG replicated proxy pilot would remain a pose-stability study, not an affinity calculation or reproduction of the unavailable 2026 input system.

The published 2026 safety conclusion also changes the research framing: further NNC2215 modeling can clarify mechanism and reproducibility, but should not be presented as discovery of a new therapeutic candidate. A new insulin design would require a separately defined molecular hypothesis and experimental validation, including evaluation of unwanted ligand interactions.

\section{Limitations}

The public structures are model-built candidates, not experimentally resolved NNC2215 structures. The 2026 PaG starting structure and trajectory are not present in the public supporting-information record checked. PDB connectivity does not define complete bond orders or partial charges. The pH 7.4 ledger is combinatorial bookkeeping without pKa calculations or measured populations. Fragment parameterization is incomplete and not independently approved. The published Figure 3d reanalysis has only three independent experiments, and PaG binding and clamp records were not refit from raw data. No receptor complex, target MD, free-energy calculation, new wet-lab experiment, or clinical outcome was generated.

These limitations preclude interpreting the present work as evidence of binding affinity, drug efficacy, safety, or clinical utility.

\section{Conclusions}

The published-data arithmetic is consistent with the rounded 2024 glucose-response summary, and a 507-heavy-atom source graph has been mapped across both publisher model-built candidates with alternatives retained. Chemical microstates and full hybrid force-field parameters remain unresolved. The current decision is no-go for target-binding MD. The next meaningful milestone is independent chemistry and parameter review followed, only if approved, by complete-system topology and energy checks. This is a model-readiness report on an existing published molecule; it does not describe a new drug.

\section{Data availability, authorship, and declarations}

The source analyses, audit scripts, model records, and machine-readable outputs are available in the project research repository at the evidence-baseline commit cited in reference [5]. The publisher model record is cited in [3], and the 2026 supporting-information record in [4]. No new human or animal experiments were conducted. A DOI-minted archive and journal-specific manuscript formatting remain to be prepared before submission.

Author list, affiliations, contributions, funding/support, and competing-interest disclosures have not been confirmed and must be completed by the contributors before journal submission. This is an internal review draft, not a submitted or accepted article.

\newpage
\begin{thebibliography}{9}
\small

\bibitem{1} Hoeg-Jensen T, Kruse T, Brand CL, et al. Glucose-sensitive insulin with attenuation of hypoglycaemia. Nature. 2024;634:944-951. doi:10.1038/s41586-024-08042-3. \href{https://doi.org/10.1038/s41586-024-08042-3}{Link}.

\bibitem{2} Sturis J, Alifrangis L, Nielsen PK, et al. Interaction between Paracetamol Glucuronide and a Glucose-Sensitive Insulin Molecule (NNC2215): Increased Risk of Hypoglycemia. ACS Pharmacol Transl Sci. 2026;9(6):1497-1504. doi:10.1021/acsptsci.6c00069. \href{https://doi.org/10.1021/acsptsci.6c00069}{Link}.

\bibitem{3} Johansson E, Granata D. 3D models of NNC2215. Springer Nature Figshare. 2024. doi:10.6084/m9.figshare.26526460.v1. \href{https://doi.org/10.6084/m9.figshare.26526460.v1}{Link}.

\bibitem{4} Sturis J, et al. Supporting Information: Interaction between Paracetamol Glucuronide and a Glucose-Sensitive Insulin Molecule (NNC2215): Increased Risk of Hypoglycemia. ACS Figshare. 2026. doi:10.1021/acsptsci.6c00069.s001. \href{https://doi.org/10.1021/acsptsci.6c00069.s001}{Link}.

\bibitem{5} Dia Health One research repository. Evidence baseline main commit 9675159eefc60bc4a8a12661bce208b99d342246. 2026. \href{https://github.com/pallagip/dia-health-one/tree/9675159eefc60bc4a8a12661bce208b99d342246/research}{Link}.

\bibitem{6} Wang J, Wolf RM, Caldwell JW, Kollman PA, Case DA. Development and testing of a general amber force field. J Comput Chem. 2004;25(9):1157-1174. doi:10.1002/jcc.20035. \href{https://doi.org/10.1002/jcc.20035}{Link}.

\bibitem{7} Jakalian A, Jack DB, Bayly CI. Fast, efficient generation of high-quality atomic charges. AM1-BCC model: II. Parameterization and validation. J Comput Chem. 2002;23(16):1623-1641. doi:10.1002/jcc.10128. \href{https://doi.org/10.1002/jcc.10128}{Link}.

\end{thebibliography}

\end{document}
